\section{Related Work}

\subsection{Test-Time Adaptation}
TTA adapts a trained model using unlabeled target samples available only during
inference. A simple and competitive approach recomputes batch-normalization
statistics from corrupted test batches~\cite{schneider2020corruptions}.
Tent extends this principle by minimizing prediction entropy while updating
normalization affine parameters online~\cite{wang2021tent}. Such objectives are
effective under covariate shift, but confident errors can be amplified when the
stream contains unknown classes. The corruption benchmarks used in this work
derive from CIFAR-10-C and CIFAR-100-C, which evaluate robustness to common,
non-adversarial corruptions~\cite{hendrycks2019benchmark}.

\subsection{Open-World TTA and STAMP}
Open-world TTA couples adaptation on known classes with detection of samples
that should not participate in the update. STAMP addresses this setting using
stable memory replay~\cite{yu2024stamp}. It aggregates predictions across
augmentations, filters inconsistent or high-entropy examples, balances the
memory across pseudo-classes, and minimizes a confidence-weighted entropy loss
on replayed samples. Our work does not redesign this adaptation pipeline.
Instead, it isolates the admission score, asking whether a detector developed
outside TTA transfers successfully when the entire known-class distribution is
also corrupted.

\subsection{Out-of-Distribution Scoring}
OOD detection provides several alternatives to entropy or maximum softmax
probability. OpenMax uses extreme-value modeling of activation-space distances
to assign probability to an unknown class~\cite{bendale2016openmax}. Deep KNN
detects OOD samples by their non-parametric distance from in-distribution
features and avoids an explicit Gaussian assumption~\cite{sun2022knn}. This
motivated our first intervention, but its usual source-reference geometry may
become unreliable when corruption shifts all test features.

Energy-based OOD detection instead derives a scalar directly from classifier
logits and was proposed to mitigate the overconfidence of softmax-based
scores~\cite{liu2020energy}. Since it requires neither a stored reference bank
nor an additional forward pass, it forms our second intervention. We also
retain two optional mechanisms for analysis: a Weibull-tail threshold inspired
by OpenMax~\cite{bendale2016openmax}, and evidential vacuity based on a
Dirichlet representation of class probabilities~\cite{sensoy2018evidential}.
Neither optional mechanism is included in the first-stage results.

\subsection{Conformal Calibration Under Shift}
Conformal methods calibrate predictive decisions from observed errors, but
standard conformal guarantees rely on an exchangeability assumption that is
inappropriate for a changing test stream. Adaptive conformal inference (ACI)
updates its calibration parameter online and targets long-run coverage under an
unknown, time-varying data-generating process~\cite{gibbs2021aci}. Our third
intervention adapts this controller to memory admission: rather than proposing
another OOD score, it adjusts the score quantile to target a specified
outlier-rejection rate. This use is experimental; the original ACI guarantees
concern conformal prediction coverage, not STAMP's OOD-detection metrics.
